{ % Adapted from RU 53535f9, section2_6.tex; editorial integration 2026-10-24.
\subsection{Profiling parallel programs}
\label{profiling:section}
\par Profiling collects information about a program's execution: time spent in routines, branch-prediction events, cache misses and related measurements. A profiler helps determine where the program spends resources and which parts prevent useful scaling.
\par\textbf{The historical Intel Parallel Amplifier workflow.} The Russian source describes three complementary views:
\begin{enumerate}
\item Hotspot analysis identifies heavily used code and examines its call stacks.
\item Concurrency analysis examines how effectively the program uses parallel execution.
\item Lock and wait analysis locates time spent waiting for synchronization rather than doing useful work.
\end{enumerate}
\par Together, these views help explain processor utilization and guide the next optimization. Change one cause at a time, check numerical correctness, and repeat the measurements under a stated workload. A profiler's own overhead and access to hardware counters must also be considered when interpreting results.
\par\textbf{Editorial note, 2026.} Intel Parallel Amplifier is the name used in the archived text, not a requirement for new experiments. Intel's current tool is VTune Profiler, with hotspot, memory and threading analysis: \url{https://www.intel.com/content/www/us/en/developer/tools/oneapi/vtune-profiler.html}. Tool support depends on the operating system, processor and analysis mode. The browser laboratory is not a replacement for native hardware-counter profiling.
\par Historical reading retained from the source: \url{https://www.ixbt.com/soft/intel-parallel-amplifier.shtml}.
}
